\section*{第 \ref{ch:g9:microbes-and-infection} 章——微生物与感染}

\begin{solution}{exo:g9:microbes-and-infection:1}
\omterm{def:g9:microbes-and-infection:sorted}{细菌}（无\omterm{def:g6:cell-unit-of-life:parts}{细胞核}的\omterm{def:g6:cell-unit-of-life:unicellular}{单细胞}）、\omterm{def:g6:cell-unit-of-life:unicellular}{单细胞}\omterm{def:g6:decomposers-and-soil:decomposers}{真菌}
如\omterm{ex:g6:producing-our-food:bread}{酵母}和池塘中的其他\omterm{def:g6:cell-unit-of-life:unicellular}{单细胞}生命，还有
\omterm{def:g9:microbes-and-infection:sorted}{病毒}。两者性质不同：\omterm{def:g9:microbes-and-infection:sorted}{病毒}没有\omterm{def:g5:first-look-at-cells:cell}{细胞}，本身没有生命，
只是一段包裹起来的\omterm{def:g9:heredity-and-traits:heredity}{遗传}信息，借用活\omterm{def:g5:first-look-at-cells:cell}{细胞}的
机器来工作。
\end{solution}

\begin{solution}{exo:g9:microbes-and-infection:2}
\omterm{def:g1:the-five-senses:sense}{皮肤}和肠道中定居着数以万亿计的无害\omterm{def:g9:microbes-and-infection:sorted}{细菌}。
它们的作用包括：帮助\omterm{def:g4:journey-of-food:digestion}{消化}，并占据
生存空间——抢占位点，挤走\omterm{prop:g9:microbes-and-infection:mostly}{病原体}。
\end{solution}

\begin{solution}{exo:g9:microbes-and-infection:3}
传播（咳嗽飞沫）；侵入（伤口、\omterm{def:g7:what-is-respiration:respiration}{呼吸}道）；
增殖（\omterm{def:g9:microbes-and-infection:sorted}{细菌}每小时分裂一次、数量倍增，
或被劫持的\omterm{def:g5:first-look-at-cells:cell}{细胞}不断复制）；疾病（化脓手指的
搏动痛、感冒引起的喉咙痛）。
\end{solution}

\begin{solution}{exo:g9:microbes-and-infection:4}
伤口中的\omterm{def:g9:microbes-and-infection:sorted}{细菌}在温暖潮湿的组织\omterm{def:g5:first-look-at-cells:cell}{细胞}之间增殖，
并毒害周围组织。感冒\omterm{def:g9:microbes-and-infection:sorted}{病毒}
在\omterm{def:g7:what-is-respiration:respiration}{呼吸}道内壁\omterm{def:g5:first-look-at-cells:cell}{细胞}内增殖，
每个被劫持的\omterm{def:g5:first-look-at-cells:cell}{细胞}都会释放出大量拷贝。
\end{solution}

\begin{solution}{exo:g9:microbes-and-infection:5}
\omterm{def:g1:the-five-senses:sense}{皮肤}：侵入时的屏障。黏膜的“扫帚”与“冲洗”
（\omterm{def:g7:what-is-respiration:respiration}{呼吸}道清扫、泪液、胃酸）：侵入与早期
增殖。\omterm{prop:g9:microbes-and-infection:mostly}{常驻菌群}：侵入——占据
生存空间。\omterm{def:g1:healthy-habits:hygiene}{卫生}措施（肥皂、烹饪、冷藏）：传播
与增殖。
\end{solution}

\begin{solution}{exo:g9:microbes-and-infection:6}
传播——产褥热的\omterm{prop:g9:microbes-and-infection:mostly}{病原体}随医生的手，
从解剖室被带到产房。巴斯德证明了
这些“旅行者”本身：\omterm{def:g9:microbes-and-infection:sorted}{微生物}是腐败、疾病和
\omterm{def:g9:microbes-and-infection:stages}{感染}的起因，而加热则是阻止它们的
克星。
\end{solution}

\begin{solution}{exo:g9:microbes-and-infection:7}
抗生素使\omterm{def:g9:microbes-and-infection:sorted}{细菌}“饥饿”并破坏\omterm{def:g9:microbes-and-infection:sorted}{细菌}——\omterm{def:g9:microbes-and-infection:sorted}{细菌}是独立的
\omterm{def:g5:first-look-at-cells:cell}{细胞}，拥有自己的机制可供毒害。感冒\omterm{def:g9:microbes-and-infection:sorted}{病毒}
没有自己的机制，它依赖\emph{你的}\omterm{def:g5:first-look-at-cells:cell}{细胞}运作，
针对其运作机制的药物也会作用于你的\omterm{def:g5:first-look-at-cells:cell}{细胞}——
这也是为什么治疗\omterm{def:g9:microbes-and-infection:sorted}{细菌}的药物对\omterm{def:g9:microbes-and-infection:sorted}{病毒}
毫无作用。
\end{solution}

\begin{solution}{exo:g9:microbes-and-infection:8}
饭前洗手：阻断手到口的传播。
如厕后洗手：阻断肠道途径中“排出—再进入”的循环。
咳嗽时捂住口鼻（打喷嚏时用肘部）：从源头阻断
飞沫传播。（刷牙和饮食习惯可阻止
口腔中的增殖。）
\end{solution}

\begin{solution}{exo:g9:microbes-and-infection:9}
\omterm{prop:g9:microbes-and-infection:mostly}{病原体}：鸡肉上的食物中毒\omterm{def:g9:microbes-and-infection:sorted}{细菌}。
传播：食物本身。侵入：肠道。增殖：
三小时温热车程——适宜\omterm{def:g6:exploring-our-environment:conditions}{条件}下\omterm{def:g9:microbes-and-infection:sorted}{细菌}每
半小时左右倍增一次。最便宜的阻断方法：
冷链——使用保温箱（或新鲜食用）；其次是
食用前彻底加热。
\end{solution}

\begin{solution}{exo:g9:microbes-and-infection:10}
因为手是在不同“隔间”之间传递\omterm{prop:g9:microbes-and-infection:mostly}{病原体}的“信使”：
塞麦尔维斯的医生按任何日常\omterm{def:g3:sorting-living-things:criterion}{标准}都算干净，
却把病房的致命\omterm{def:g1:healthy-habits:germs}{病菌}从尸体带到了产妇身上。真正
关键的是\emph{在边界处}的那一次洗手——生鸡肉与
沙拉之间、厕所与餐桌之间，那里传播途径
否则就会接通。
\end{solution}

\begin{solution}{exo:g9:microbes-and-infection:11}
抗生素攻击的是增殖——破坏正在分裂的\omterm{def:g9:microbes-and-infection:sorted}{细菌}。
若使用不规范，就成了一次中途进行的选择
检验：敏感菌死亡，耐药的少数留存下来，
并占据病房，培育出下一位患者
所遇到的\omterm{prop:g9:microbes-and-infection:mostly}{病原体}。
\end{solution}

\begin{solution}{exo:g9:microbes-and-infection:12}
\omterm{def:g9:microbes-and-infection:sorted}{病毒}的“文本”是一段\omterm{def:g9:heredity-and-traits:heredity}{遗传}密码，与每个\omterm{def:g5:first-look-at-cells:cell}{细胞}
所用的字母相同（见\cref{rem:g9:genes-and-dna:universal}）。“印刷机”则是
活\omterm{def:g5:first-look-at-cells:cell}{细胞}的读写与合成机制，被侵入的文本驱动它
印出自身的拷贝。
生命的边界：单独存在时，\omterm{def:g9:microbes-and-infection:sorted}{病毒}在
\cref{prop:g1:living-or-not:signs}所列的特征中一项也不具备——不摄食、
不生长、不分裂；借用\omterm{def:g5:first-look-at-cells:cell}{细胞}时，它演出的是一场
繁殖的“拙劣模仿”。难以用药：\omterm{def:g9:microbes-and-infection:sorted}{病毒}自身没有
可供毒害的机制，唯一的“印刷机”是你的\omterm{def:g5:first-look-at-cells:cell}{细胞}。
\end{solution}

\begin{solution}{pb:g9:microbes-and-infection:1}
\textbf{1.} 飞沫传播——咳嗽和说话时产生的飞沫，
并借助乒乓球拍等共用表面传播；
通过\omterm{def:g7:what-is-respiration:respiration}{呼吸}道侵入。

\textbf{2.} 增殖：每个病例中被劫持的\omterm{def:g7:what-is-respiration:respiration}{呼吸}道内壁
\omterm{def:g5:first-look-at-cells:cell}{细胞}会印出拷贝，进而\omterm{def:g9:microbes-and-infection:stages}{感染}下一批病例——
这种倍增是\omterm{def:g9:microbes-and-infection:sorted}{病毒}借用的机制在逐级累积。

\textbf{3.} 疾病：来自被劫持、功能衰竭的内壁\omterm{def:g5:first-look-at-cells:cell}{细胞}受损——
并在一定程度上来自防御系统自身的战斗，
下一章将详细说明。

\textbf{4.} 由于集会使每个班级在一小时内的飞沫范围相互连通，
传播程度随人群拥挤而增加；
学校最密集的两次混合，
正是各区域之间的传播途径。

\textbf{5.} 患病学生回家：在源头
切断传播。开窗：稀释飞沫剂量——切断传播。
勤洗手：切断经表面接触的途径——切断传播。
擦拭器材：切断共用物品的途径——
切断传播。（这些措施都针对同一个成本最低的阶段。）

\textbf{6.} “抗生素杀灭\omterm{def:g9:microbes-and-infection:sorted}{细菌}；这种\omterm{prop:g9:microbes-and-infection:mostly}{病原体}是
寄生在孩子自身\omterm{def:g5:first-look-at-cells:cell}{细胞}上的\omterm{def:g9:microbes-and-infection:sorted}{病毒}——这种药物
无法作用于它；即使如此仍要撒药，也只会
对每个人的常驻\omterm{def:g9:microbes-and-infection:sorted}{细菌}进行一轮选择筛选，
为将来真正需要抗生素的那一天培育
抗性。”

\textbf{7.} \omterm{prop:g9:microbes-and-infection:mostly}{病原体}：食物中毒\omterm{def:g9:microbes-and-infection:sorted}{细菌}。传播途径：
甜点。侵入：肠道。增殖：未冷藏的
数小时——\omterm{def:g6:exploring-our-environment:conditions}{温度}、\omterm{def:g6:exploring-our-environment:conditions}{湿度}、食物与时间，\omterm{def:g6:exploring-our-environment:conditions}{条件}样样具备。
阻断：冰箱——拒绝提供\omterm{def:g6:exploring-our-environment:conditions}{条件}；两次
暴发仅共享同一所学校，别无其他共同点。

\textbf{8.} 疫情图是把传播过程可视化：
谁坐在哪里便勾勒出路径——座位排、同桌、
集会——路径一经画出，便显示该在何处
切断。随机治疗只是对付病例；疫情图则是对付
传播本身。

\textbf{9.} 传播：首例病例的飞沫，先经座位排和同桌，
再经集会传至全校。侵入：
经\omterm{def:g7:what-is-respiration:respiration}{呼吸}道，逐班进行。增殖：第一周内每两天
翻一番。疾病与消退：发热高峰出现，
防控措施切断传播途径，易感人群
逐渐减少——疫情曲线开始下降。

\textbf{10.} 防控措施的阻断使传播途径断绝——
而易感人群本身也比名册上的人数少：学生
因去年冬天\omterm{def:g9:microbes-and-infection:stages}{感染}过近缘毒株而具有免疫力，
无法被点燃；可被\omterm{def:g9:microbes-and-infection:stages}{感染}者太少时，疫情便会
消亡——这是下一章关于免疫记忆的预告。

\textbf{11.} 塞麦尔维斯的做法：切断传播途径——
居家隔离、洗手、擦拭球拍。巴斯德的做法：
认清传播者——明确\omterm{prop:g9:microbes-and-infection:mostly}{病原体}，拒绝为它
提供\omterm{def:g6:exploring-our-environment:conditions}{条件}，使对策针对其种类，而不是沿用
惯例。

\textbf{12.} “疫情在传播阶段止住成本最低——
肥皂、通风、距离与封闭的
冷藏箱，胜过任何等到
增殖阶段才出手的药物。”
\end{solution}
